\BourbakiBackSection{历史注记}

（方括号中的编号指本注记末尾的参考文献。）

关于一致空间的主要概念与结果是逐渐从实变量理论中产生的，只是到了近年它们才得到系统的研究。Cauchy 在为级数理论奠定严格基础的尝试中（参看第一章与第四章的历史注记），以一条他似乎认为是自明的原则作为出发点，即：序列 \((a_{n})\) 收敛的充分必要条件是，只要 n 充分大，\(|a_{n+p} - a_{n}|\) 就可任意小（例如见 {[}2{]}）。他与 Bolzano {[}1{]} 无疑是最早明确陈述这条原则并认识到其重要性的人；因此，满足所述条件的实数序列被称为“Cauchy 序列”，并按推广的含义，这一名称也用于度量空间（第九章）中满足下列条件的点列 \((x_{n})\)：只要 n 充分大，\(x_{n+p}\) 到 \(x_{n}\) 的距离就可任意小；最终，本章所研究的 Cauchy 序列的推广也获得了“Cauchy 滤子”这一名称。

后来，当实数的直观概念不再被认为足够，人们为了给分析提供坚实的基础而试图用有理数来定义实数时，正是 Cauchy 的原则提供了十九世纪下半叶提出的各种定义中最富成果的一个。这个定义属于 Cantor {[}3{]}（Heine {[}5{]} 沿 Cantor 思想的路线以及 Méray 等人又各自独立地发展了它），其内容是：使每个有理数的 Cauchy 序列（Cantor 术语中的“基本序列”）对应于一个实数；两个有理数的 Cauchy 序列 \((a_{n})\) 与 \((b_{n})\) 对应于同一个实数，当且仅当 \(|a_{n} - b_{n}|\) 趋于零。这里的本质思想是：从某种观点（事实上，即本章 § 1 第 1 目例 1 中定义的“一致结构”的观点）看来，有理数集 Q 是“不完备”的，而实数集是由 Q 经“完备化”得到的“完备”集。

另一方面，Heine 在很大程度上受 Weierstrass 与 Cantor 思想的启发，第一个定义了一个或多个实变量的实值函数的一致连续性 {[}4{]}，并证明了在 R 的有界闭区间上连续的每个实值函数都在该区间上一致连续：这就是“Heine 定理”。由 § 4 定理 2，这一结果与 R 的有界闭区间的紧性（“Borel-Lebesgue 定理”，第四章 § 2 定理 2；参看第一章与第四章的历史注记）相关，而且 Heine 对其定理的证明稍加修改也可用来证明 Borel-Lebesgue 定理（这在某些作者看来已足以成为把后一定理称为“Heine-Borel 定理”的理由）。

当度量空间开始被研究时——先是特殊情形，然后是一般情形（参看第九章）——这些思想便被推广到更一般的空间；在度量空间中给定了距离（即满足某些公理的点对的实值函数），它同时定义了一个拓扑和一个一致结构。Fréchet 第一个给出了这些空间的一般定义，他认识到了 Cauchy 原则的重要性 {[}6{]}，并对度量空间证明了与 § 4 定理 3 等价的一个定理（{[}6{]} 与 {[}7{]}）。Hausdorff 在他的《Mengenlehre》（{[}8{]}；并参看 {[}8 a{]}）中相当可观地发展了度量空间的理论，并且特别地认识到，可以把前面描述的 Cantor 构造应用于这些空间，从而从一个“非完备”度量空间（即 Cauchy 原则在其中不成立的空间）得到一个“完备”度量空间。

度量空间是“一致空间”的一种特殊类型；后者是在相当晚近的时候才由 A. Weil 在完全一般的形式下定义的 {[}9{]}。在此之前，关于“一致结构”的概念和结果只能结合度量空间来使用，这一事实说明了度量空间与可度量化空间（特别是紧可度量化空间）在许多现代拓扑著作中所起的重要作用，尽管在这些问题中距离并没有真正的用处。一旦有了一致空间的定义，把这些空间上几乎整个（例如 Hausdorff 所给出的）度量空间理论加以推广就没有什么困难了（特别当手头还具有滤子概念时更是如此；类似地，例如可以把 Alexandroff-Hopf 的《Topologie》 {[}10{]} 中关于紧度量空间的结果推广到任意紧空间）。我们在本章所做的正是这件事；特别地，关于一致空间完备化的定理（§ 3 定理 3）不过是 Cantor 实数构造的不带任何本质修改的移植。

